\documentclass[12pt]{article}
\usepackage[a3paper,landscape]{geometry}
\usepackage{lipsum}
\usepackage{lmodern}
\usepackage{enumerate}
\usepackage[poster]{tcolorbox}
\tcbuselibrary{listings} % <- replace by \tcbuselibrary{listings}, if minted does not work for you
\pagestyle{empty}
\begin{document}
	\begin{tcbposter}[
		coverage = {
			spread,
			interior style={top color=white,bottom color=gray!50!}
		},poster = {showframe=true,columns=6,rows=6},
			boxes = {
			enhanced standard jigsaw,sharp corners=downhill,arc=1mm,boxrule=0.25mm,
			colback=white,opacityback=1,colframe=black,
			title style={color=black},
			fonttitle=\bfseries\Large\scshape
		}
				]
		% Ab hier füge ich die Poster Inhalte ein 
	%...
	\posterbox{name=title,column=2,span=4,below=top}{
		\resizebox{12cm}{!}{\bfseries\Huge Cheat Sheet - PLM}\\[1mm]
		Marc.Weidner@protonmail.ch
	}
	%----------------- Fußzeile 
	\posterbox[adjusted title=Datum]{name=datum,column=1,span=1,above=bottom}{}
	\posterbox[adjusted title=References]{name=references,column=2,span=1,above=bottom}{}
	\posterbox[adjusted title=Erzeugt]{name=erzeugt,column=3,span=2,above=bottom}{}
	\posterbox[adjusted title=Klassifizierung]{name=class,column=5,span=1,above=bottom}{}
	\posterbox[adjusted title=Seite]{name=page,column=6,span=1,above=bottom}{\pagenumbering{1/1}}
	%--------------- BOX lInks oben
	\posterbox[adjusted title=Beschreibung]{name=desc,
		sequence=1 between top and datum then 2 between title and references}{
		See [\ref{litA}]: \lipsum[1]
		\begin{center}
			\tikz \draw[thick,rounded corners=8pt]
			(0,0)--(0,2)--(1,3.25)--(2,2)--(2,0)--(0,2)--(2,2)--(0,0)--(2,0);
			\quad by [\ref{litB}]
		\end{center}
		\lipsum[2-3]\par
		See [\ref{litC}]:
		\lipsum[4]
		\begin{center}
			\tikz \shadedraw [left color=red,right color=blue]
			(0,0) rectangle (2,2);
		\end{center}
		That's all.
	}
	%...
	%...bild rechts oben
	\posterbox[adjusted title=Central Picture,
	interior style={fill overzoom image=blueshade.png}]
	{name=picture,column=6,below=top}{}
	%..
	%--------------- BOX Center mit Graphiken
	\posterbox[adjusted title=Center Box 1,halign=center]{name=center,column=3,span=3,above=erzeugt}{
		\begin{tikzpicture}[very thick,radius=2cm]
			\begin{scope}
				\path[draw=black,fill=white] (0,0) circle;
				\path[fill=red] (0,0) -- (2,0) arc [start angle=0, end angle=30];
			\end{scope}
			\begin{scope}[xshift=5cm]
				\path[draw=black,fill=white] (0,0) circle;
				\path[fill=red] (0,0) -- (2,0) arc [start angle=0, end angle=70];
			\end{scope}
			\begin{scope}[xshift=10cm]
				\path[draw=black,fill=white] (0,0) circle;
				\path[fill=red] (0,0) -- (2,0) arc [start angle=0, end angle=110];
			\end{scope}
			\begin{scope}[xshift=15cm]
				\path[draw=black,fill=white] (0,0) circle;
				\path[fill=red] (0,0) -- (2,0) arc [start angle=0, end angle=240];
			\end{scope}
		\end{tikzpicture}
	}
	%..
	
	\end{tcbposter}
\end{document}





